% =====================================================================
%  CHAPTER 6 -- RESULTS AND DISCUSSION
%
%  RUBRIC NOTE (R8): results must be "presented with supporting figures
%  and graphics". A results chapter with no figure cannot reach the top
%  band, however good the prose is. Figure and table slots are reserved
%  below -- fill at least two of them.
% =====================================================================
\chapter{RESULTS AND DISCUSSION}
\label{ch:results}

% TODO: one orientation paragraph.
[CHAPTER OPENING: state what was evaluated, against what criteria, and
how the results are organised in this chapter.]

% ---------------------------------------------------------------------
\section{EVALUATION CRITERIA}
\label{sec:evaluation-criteria}
% ---------------------------------------------------------------------

% TODO: state what you measured and how, BEFORE showing numbers. A
% reader has to know the measurement method to trust the measurement.
[CRITERIA: the metrics used to evaluate the system, how each was
measured, the conditions under which measurements were taken, and how
many times each measurement was repeated.]

% ---------------------------------------------------------------------
\section{RESULTS}
\label{sec:results}
% ---------------------------------------------------------------------

% TODO: present the results. Numbers first, interpretation after.
[RESULTS NARRATIVE: present what the system achieved against each
criterion from Section~\ref{sec:evaluation-criteria}. Report the
measured values plainly here; save the interpretation for
Section~\ref{sec:discussion}.]

% ---- Results table ---------------------------------------------------
\begin{table}[H]
  \centering
  \caption{Summary of evaluation results}
  \label{tab:results-summary}
  \small
  \begin{tabularx}{\textwidth}{|l|X|l|l|}
    \hline
    \textbf{Metric} & \textbf{Description} & \textbf{Target} & \textbf{Achieved} \\
    \hline
    [METRIC 1] & [WHAT IT MEASURES] & [TARGET] & [RESULT] \\
    \hline
    [METRIC 2] & [WHAT IT MEASURES] & [TARGET] & [RESULT] \\
    \hline
    [METRIC 3] & [WHAT IT MEASURES] & [TARGET] & [RESULT] \\
    \hline
    % TODO: one row per metric.
  \end{tabularx}
\end{table}

% ---- Supporting figure slots (fill at least two) ---------------------
\begin{figure}[H]
  \centering
  \fypdiagram[7cm]{[INSERT RESULTS GRAPH 1 HERE -- e.g. accuracy, response time or throughput chart]}
  \caption{[Describe what this graph shows]}
  \label{fig:results-graph-1}
\end{figure}

% TODO: every figure needs a sentence in the body text that REFERS to
% it. A figure nobody points at does not earn its marks.
[Refer to the figure here, e.g.: As shown in
Figure~\ref{fig:results-graph-1}, ...]

\begin{figure}[H]
  \centering
  \fypdiagram[7cm]{[INSERT RESULTS GRAPH 2 HERE -- e.g. comparison against a baseline or against existing systems]}
  \caption{[Describe what this graph shows]}
  \label{fig:results-graph-2}
\end{figure}

[Refer to Figure~\ref{fig:results-graph-2} here.]

% ---------------------------------------------------------------------
\section{DISCUSSION}
\label{sec:discussion}
% ---------------------------------------------------------------------

% TODO: interpret the numbers. What do they mean?
[INTERPRETATION: explain what the results mean. Did the system meet the
objectives set out in Section~\ref{sec:objectives}? Answer objective by
objective.]

% TODO: compare against the literature. This is what ties the report
% back to Chapter 2 and lifts the discussion above a lab report.
[COMPARISON: compare these results against the existing approaches
reviewed in Chapter~\ref{ch:literature-review}. The most direct
comparison available to you is the hybrid recommender for partner
selection \cite{zhang2021hybrid}, which solves a related matching
problem in the same domain. Where does this system do better, where
does it do worse, and why?]

% TODO: be honest about limitations. Examiners reward this.
[LIMITATIONS: state the limitations of the system and of the evaluation
itself -- small sample size, restricted test environment, untested edge
cases. An objective account of limitations is explicitly rewarded by
the grading rubric.]
